% GENERATED FILE. Edit canonical lessons, then run npm run generate:books.
% canonical-source-hash: fnv1a64:2875dd5345a2bae4
% canonical-lessons: HI-C439-tap, HI-C439-para, HI-C439-tapman, HI-C439-digri, HI-C439-tapna

\chapter{Heat}
\label{ch:hi-heat}
\hlchaptermodality{439}

\begin{chapteropening}
\textbf{By the end of this chapter:} \emph{I can talk about the heat and the temperature in degrees, and roads scorching in the sun.}

Complete the last lesson of chapter 439: I can talk about the heat and the temperature in degrees, and roads scorching in the sun.
\end{chapteropening}

\section[\texorpdfstring{\hi{ताप} (tāp)}{ताप (tāp)}]{\hi{ताप} (tāp) — heat; a fever (formal)}

\label{lesson:HI-C439-tap}



\begin{quote}
Before the new one: say the Hindi for snowfall, then the Hindi for to wrap oneself in.
\end{quote}

\subsection*{You'll want to know: \hi{ताप}}
\textbf{\hi{ताप}} — \emph{tāp} — \textquotedblleft{}heat; a fever (formal)\textquotedblright{}.

\begin{quote}
\hi{दोपहर} \hi{का} \hi{ताप} \hi{सहना} \hi{मुश्किल} \hi{था।}
\end{quote}

\begin{quote}
\emph{dopahar kā tāp sahnā muśkil thā.}
\end{quote}

\begin{quote}
The midday heat was hard to bear.
\end{quote}

Masculine.

\begin{cousinweb}
\textbf{\hi{ताप}} (\emph{tāp}) is a tatsama, Sanskrit \emph{tāpa}, \textquotedblleft{}heat, glow\textquotedblright{} (Monier-Williams), from the root \emph{tap}, \textquotedblleft{}to be hot\textquotedblright{}, which comes back in a verb later in this chapter.
\end{cousinweb}

\begin{grammarlens}[title={What you've built}]
One more word for getting around.
\end{grammarlens}

\subsection*{Your turn}
\begin{itemize}
\raggedright
  \item \emph{Say it:} \emph{tāp}
  \item \emph{Say it:} \emph{tāp}, once more
  \item \emph{Say it:} \emph{dopahar kā tāp sahnā muśkil thā.} — The midday heat was hard to bear.
  \item \emph{Recall:} say the Hindi for snowfall, then the Hindi for to wrap oneself in, then say \emph{tāp} again
\end{itemize}

\begin{tcolorbox}[breakable,colback=teal!4,colframe=teal!35!black,title={Before you move on}]
What does \hi{ताप} mean? (\textquotedblleft{}Heat; a fever (formal)\textquotedblright{}.) Say it once more.
\end{tcolorbox}
\section[\texorpdfstring{\hi{पारा} (pārā)}{पारा (pārā)}]{\hi{पारा} (pārā) — mercury; the temperature}

\label{lesson:HI-C439-para}



\begin{quote}
Before the new one: say the Hindi for to wrap oneself in, then the Hindi for heat (formal).
\end{quote}

\subsection*{You'll want to know: \hi{पारा}}
\textbf{\hi{पारा}} — \emph{pārā} — \textquotedblleft{}mercury; the temperature\textquotedblright{}.

\begin{quote}
\hi{आज} \hi{पारा} \hi{बहुत} \hi{ऊपर} \hi{चला} \hi{गया।}
\end{quote}

\begin{quote}
\emph{āj pārā bahut ūpar calā gayā.}
\end{quote}

\begin{quote}
The temperature shot up today.
\end{quote}

Masculine.

\begin{cousinweb}
\textbf{\hi{पारा}} (\emph{pārā}) comes from Sanskrit \emph{pārada}, \textquotedblleft{}quicksilver\textquotedblright{} (Monier-Williams); the mercury in a thermometer gave the weather news its \hi{पारा} \hi{चढ़ा}, the temperature rose.
\end{cousinweb}

\begin{grammarlens}[title={What you've built}]
One more word for getting around.
\end{grammarlens}

\subsection*{Your turn}
\begin{itemize}
\raggedright
  \item \emph{Say it:} \emph{pārā}
  \item \emph{Say it:} \emph{pārā}, once more
  \item \emph{Say it:} \emph{āj pārā bahut ūpar calā gayā.} — The temperature shot up today.
  \item \emph{Recall:} say the Hindi for to wrap oneself in, then the Hindi for heat (formal), then say \emph{pārā} again
\end{itemize}

\begin{tcolorbox}[breakable,colback=teal!4,colframe=teal!35!black,title={Before you move on}]
What does \hi{पारा} mean? (\textquotedblleft{}Mercury; the temperature\textquotedblright{}.) Say it once more.
\end{tcolorbox}
\section[\texorpdfstring{\hi{तापमान} (tāpmān)}{तापमान (tāpmān)}]{\hi{तापमान} (tāpmān) — temperature}

\label{lesson:HI-C439-tapman}



\begin{quote}
Before the new one: say the Hindi for heat (formal), then the Hindi for mercury.
\end{quote}

\subsection*{You'll want to know: \hi{तापमान}}
\textbf{\hi{तापमान}} — \emph{tāpmān} — \textquotedblleft{}temperature\textquotedblright{}.

\begin{quote}
\hi{रात} \hi{को} \hi{तापमान} \hi{बहुत} \hi{नीचे} \hi{गिर} \hi{गया।}
\end{quote}

\begin{quote}
\emph{rāt ko tāpmān bahut nīce gir gayā.}
\end{quote}

\begin{quote}
At night the temperature dropped a long way.
\end{quote}

Masculine.

\begin{cousinweb}
\textbf{\hi{तापमान}} (\emph{tāpmān}) is a modern compound of Sanskrit \emph{tāpa}, \textquotedblleft{}heat\textquotedblright{}, and \emph{māna}, \textquotedblleft{}measure\textquotedblright{} (Monier-Williams for both): the measure of heat.
\end{cousinweb}

\begin{grammarlens}[title={What you've built}]
One more word for getting around.
\end{grammarlens}

\subsection*{Your turn}
\begin{itemize}
\raggedright
  \item \emph{Say it:} \emph{tāpmān}
  \item \emph{Say it:} \emph{tāpmān}, once more
  \item \emph{Say it:} \emph{rāt ko tāpmān bahut nīce gir gayā.} — At night the temperature dropped a long way.
  \item \emph{Recall:} say the Hindi for heat (formal), then the Hindi for mercury, then say \emph{tāpmān} again
\end{itemize}

\begin{tcolorbox}[breakable,colback=teal!4,colframe=teal!35!black,title={Before you move on}]
What does \hi{तापमान} mean? (\textquotedblleft{}Temperature\textquotedblright{}.) Say it once more.
\end{tcolorbox}
\section[\texorpdfstring{\hi{डिग्री} (ḍigrī)}{डिग्री (ḍigrī)}]{\hi{डिग्री} (ḍigrī) — a degree}

\label{lesson:HI-C439-digri}



\begin{quote}
Before the new one: say the Hindi for mercury, then the Hindi for temperature.
\end{quote}

\subsection*{You'll want to know: \hi{डिग्री}}
\textbf{\hi{डिग्री}} — \emph{ḍigrī} — \textquotedblleft{}a degree\textquotedblright{}.

\begin{quote}
\hi{आज} \hi{दोपहर} \hi{को} \hi{चालीस} \hi{डिग्री} \hi{गर्मी} \hi{थी।}
\end{quote}

\begin{quote}
\emph{āj dopahar ko cālīs ḍigrī garmī thī.}
\end{quote}

\begin{quote}
It was forty degrees this afternoon.
\end{quote}

Feminine.

\begin{cousinweb}
\textbf{\hi{डिग्री}} (\emph{ḍigrī}) is borrowed from English \emph{degree}, of heat and also of a university; it is feminine.
\end{cousinweb}

\begin{grammarlens}[title={What you've built}]
One more word for getting around.
\end{grammarlens}

\subsection*{Your turn}
\begin{itemize}
\raggedright
  \item \emph{Say it:} \emph{ḍigrī}
  \item \emph{Say it:} \emph{ḍigrī}, once more
  \item \emph{Say it:} \emph{āj dopahar ko cālīs ḍigrī garmī thī.} — It was forty degrees this afternoon.
  \item \emph{Recall:} say the Hindi for mercury, then the Hindi for temperature, then say \emph{ḍigrī} again
\end{itemize}

\begin{tcolorbox}[breakable,colback=teal!4,colframe=teal!35!black,title={Before you move on}]
What does \hi{डिग्री} mean? (\textquotedblleft{}A degree\textquotedblright{}.) Say it once more.
\end{tcolorbox}
\section[\texorpdfstring{\hi{तपना} (tapnā)}{तपना (tapnā)}]{\hi{तपना} (tapnā) — to scorch; to be burning hot}

\label{lesson:HI-C439-tapna}



\begin{quote}
Before the new one: say the Hindi for temperature, then the Hindi for a degree.
\end{quote}

\subsection*{You'll want to know: \hi{तपना}}
\textbf{\hi{तपना}} — \emph{tapnā} — \textquotedblleft{}to scorch; to be burning hot\textquotedblright{}.

\begin{quote}
\hi{जून} \hi{में} \hi{सड़कें} \hi{धूप} \hi{में} \hi{तपती} \hi{हैं।}
\end{quote}

\begin{quote}
\emph{jūn meṁ saṛkeṁ dhūp meṁ taptī haiṁ.}
\end{quote}

\begin{quote}
In June the roads bake in the sun.
\end{quote}

Past (it …): \hi{वह} \hi{छत} \hi{धूप} \hi{में} \hi{तपी}, \emph{vah chat dhūp meṁ tapī}.

\begin{cousinweb}
\textbf{\hi{तपना}} (\emph{tapnā}) is a tadbhava of Sanskrit \emph{tapyati}, \textquotedblleft{}is hot\textquotedblright{} (CDIAL 5684), the same root as \hi{ताप}.
\end{cousinweb}

\begin{grammarlens}[title={What you've built}]
One more word for getting around.
\end{grammarlens}

\subsection*{Your turn}
\begin{itemize}
\raggedright
  \item \emph{Say it:} \emph{tapnā}
  \item \emph{Say it:} \emph{tapnā}, once more
  \item \emph{Say it:} \emph{jūn meṁ saṛkeṁ dhūp meṁ taptī haiṁ.} — In June the roads bake in the sun.
  \item \emph{Recall:} say the Hindi for temperature, then the Hindi for a degree, then say \emph{tapnā} again
\end{itemize}

\begin{tcolorbox}[breakable,colback=teal!4,colframe=teal!35!black,title={Before you move on}]
What does \hi{तपना} mean? (\textquotedblleft{}To scorch; to be burning hot\textquotedblright{}.) Say it once more.
\end{tcolorbox}
